% HAND-WRITTEN fragment. The procedure. Timing arithmetic: one conversion is
% 1.15 us (note_adc_intro), so 4096 codes x 16 averages = 75 ms per channel and
% a 256-average zero is 0.3 ms. Flash ownership and the boot-time staging model
% come from MULTICORE-intro-castalia-2026-07.tex (hart 0 is the management and
% boot hart, the only one on the SPI-flash path, and it stages each channel
% hart's image before igniting it).

\subsubsection{Procedure}\label{sss:calibration-procedure}

Every correction is digital. No analog trim is proposed beyond the bias code the
generator already carries: the offsets are static, a stored constant removes
them to the converter's own resolution, and an analog null could do no better
without a finer converter. The one case where that reasoning fails is a
saturated channel, and it is a range condition rather than an offset one
(Section~\ref{sss:calibration-residual}).

\paragraph{Once per chip, at production test.} Steps 1 to 10, in this order,
because each one is measured through the settings the previous one fixed.

\begin{enumerate}[leftmargin=*, itemsep=2pt]
\item \textbf{Bias trim.} Electrodes isolated, all channels enabled. Step the
  6-bit bias code and read the total analog supply current; keep the code
  nearest the design current. The trim spans \SI{39.9}{\micro\ampere}
  (Table~\ref{tab:biasgen-trim}) against a \SI{12.5}{\percent} spread, so
  resolution and not range is the limit: \SI{2.7}{\percent} per code,
  \(\pm1.4\,\%\) after rounding. This is first because it moves every
  transconductance behind every constant that follows.
\item \textbf{Absolute anchor, voltage.} Route the channel's reference DAC
  output to the test pad; measure codes 0, 2048 and 4095 with a voltmeter.
  Convert the same three codes through the multiplexer in the same pass. This
  fixes the DAC's step size in volts and, with it, the converter's --- the only
  two absolute scales on the chip.
\item \textbf{Converter zero.} Multiplex the common-mode rail into the
  converter and average. The resulting code is the origin of the current axis
  and carries the converter's offset and the common-mode value together, which
  is exactly the combination the current reading needs.
\item \textbf{A1 offset.} Open the reference-electrode switch, close the
  loopback, multiplex the counter-electrode node in, and read at three
  reference codes across the working window. The intercept is the amplifier's
  input offset, the slope its finite-loop-gain tracking error. Repeat per
  channel.
\item \textbf{Current zero.} Open the loopback and the working-electrode switch
  so that the cell current is zero by construction, and read the transimpedance
  output at each gain code the channel will use. Subtracting step 3 gives the
  zero-current offset in codes.
\item \textbf{Absolute anchor, current.} Close the working-electrode switch;
  force \(\pm\) full-scale current into the pad at each stored gain code and
  read the converter. The slope is one end-to-end constant per gain code in
  amperes per code, and it absorbs the feedback resistor's \(\pm35\,\%\) sheet
  tolerance, the closed-switch stack, the transimpedance loop's finite-gain
  slope error and the converter's \(1.535\times\) gain anomaly in a single
  number. Nothing downstream needs to know how that number divides between them.
\item \textbf{On-chip gain standard.} Isolate the electrodes, close the loopback
  and the gain-reference switch, and repeat the gain measurement against the
  internal resistor at the same codes. This ties the on-chip ratio to the
  external ampere once, so that every later gain check needs no equipment.
\item \textbf{Reference DAC.} Sweep all 4096 codes into the converter and fit
  \(V = V_0 - R_{\mathrm{out}} I_{\mathrm{ref}}(\mathrm{code})\), the
  ideal-ladder reference current of Section~\ref{sss:dacr2r12-supply}. One
  stored parameter removes \SI{94}{\percent} of the \SI{25.8}{\milli\volt}
  reference movement. Do not keep the per-code residual table: measured through
  this converter each of its entries carries the same \SI{0.65}{\milli\volt} of
  quantisation as the error it corrects, so it moves the potential residual from
  \SI{1.38}{\milli\volt} rms to \SI{1.31}{\milli\volt} and costs
  4\,KiB per channel. Run the sweep regardless --- it is
  \SI{75}{\milli\second} per channel at sixteen averages, and it is where a
  broken ladder shows up.
\item \textbf{Switch resistance.} Force full-scale current with the
  working-electrode switch closed and measure the pad against the common-mode
  rail. Store the resistance; the firmware correction needs it, not a small
  value.
\item \textbf{Commit.} Write the constant block of
  Table~\ref{tab:calibration-constants} to flash with a checksum.
\end{enumerate}

\paragraph{At every power-up, on-die, no equipment.} The management hart owns
the flash path and stages each channel hart before starting it, so the constants
arrive with the image.

\begin{enumerate}[leftmargin=*, itemsep=2pt, resume]
\item Load and verify the constant block; apply the bias code before anything
  else is measured.
\item With the electrodes isolated, repeat steps 3, 4 and 5 --- the converter
  zero, A1's offset, and the current zero at the active gain code. These replace
  their stored values. Each is a difference of two averaged conversions and
  costs under a millisecond per channel.
\item Repeat step 7 against the on-chip gain resistor and scale the stored
  absolute gain by the change in the ratio. The absolute anchors from steps 2
  and 6 are not re-measurable in the field and are carried forward unchanged.
\item Compare each refreshed value against the value stored at test and raise
  the site's status flag if any has moved by more than its limit. A drifted zero
  is a usable channel; a drifted gain ratio is a damaged one.
\item Close the electrode switches. The channel is ready.
\end{enumerate}

\paragraph{During a measurement.} Re-zero between potential steps by
multiplexing the common-mode rail in for one conversion --- \SI{1.15}{\micro
\second}, against settling times three orders of magnitude longer --- and apply
the corrections in this order: subtract the zero code, multiply by the gain
constant for the active gain code, then subtract the potential lost across the
working-electrode switch and, where an impedance sweep has measured it, the
uncompensated electrolyte resistance. Potential-axis corrections apply where the
waveform is built, not per sample: the DAC code for a commanded potential comes
from the corrected transfer, and cyclic-voltammetry and pulse waveforms are
tabulated before the sweep starts.
